\documentclass[10pt,letterpaper]{article}
\usepackage[spanish,es-nodecimaldot]{babel}
\usepackage{fontspec}
\usepackage[letterpaper,top=1.8cm,bottom=1.8cm,left=1.8cm,right=1.8cm]{geometry}
\usepackage{xcolor,graphicx,float,booktabs,longtable,array,fancyhdr,hyperref,listings,lastpage}
\definecolor{accent}{HTML}{A6192E}\definecolor{muted}{HTML}{667085}\definecolor{codebg}{HTML}{F5F6F8}
\hypersetup{colorlinks=true,linkcolor=accent,urlcolor=accent,pdftitle={Manual tecnico - Comercial La Estrella},pdfauthor={Alex Ricardo Castaneda Rodriguez}}
\pagestyle{fancy}\fancyhf{}\fancyhead[L]{\color{accent}Base de Datos 1 · Proyecto 1}\fancyhead[R]{Manual técnico}\fancyfoot[L]{Alex Ricardo Castañeda Rodríguez}\fancyfoot[R]{\thepage/\pageref*{LastPage}}
\setlength{\parindent}{0pt}\setlength{\parskip}{5pt}\renewcommand{\arraystretch}{1.14}
\lstset{language=SQL,basicstyle=\ttfamily\footnotesize,backgroundcolor=\color{codebg},frame=single,breaklines=true,showstringspaces=false,keywordstyle=\bfseries\color{accent},commentstyle=\itshape\color{muted}}
\newcommand{\modelo}[2]{\begin{figure}[H]\centering\includegraphics[width=\textwidth,height=.62\textheight,keepaspectratio]{#1}\caption{#2}\end{figure}}
\begin{document}
\begin{titlepage}\thispagestyle{empty}
\noindent\color{accent}\rule{\textwidth}{8pt}\par\vspace{2cm}
\begin{center}{\Large UNIVERSIDAD DE SAN CARLOS DE GUATEMALA}\par\medskip{\large Facultad de Ingeniería · Escuela de Ciencias y Sistemas}\par\vfill
{\Huge\bfseries MANUAL TÉCNICO}\par\bigskip{\LARGE\color{accent}Diseño e implementación de una base de datos para el control de ventas}\par\medskip{\Large Comercial La Estrella}\par\vfill
\begin{tabular}{rl}\textbf{Estudiante:}&Alex Ricardo Castañeda Rodríguez\\\textbf{Carné:}&202300476\\\textbf{Curso:}&Bases de Datos 1\\\textbf{Proyecto:}&1\\\textbf{Fecha:}&10 de septiembre de 2026\end{tabular}\par\vfill
{\small Guatemala, 2026}\end{center}\noindent\color{accent}\rule{\textwidth}{4pt}
\end{titlepage}
\tableofcontents\clearpage

\section{Propósito y arquitectura}
La solución administra ubicaciones, tiendas, personas, empleados, clientes, productos, catálogo por tienda, ventas, detalles y pagos. Se implementan 19 tablas Oracle y se cargan exactamente 9,068 registros provenientes de las 19 hojas del libro oficial. Los identificadores del archivo se preservan; no se usan columnas \texttt{IDENTITY}.

La arquitectura se divide en cuatro áreas: catálogos y ubicación; personas y establecimientos; producto e inventario; transacciones y pagos. \texttt{persona} concentra los atributos compartidos por empleados y clientes. \texttt{catalogo\_producto} resuelve la relación muchos a muchos tienda--producto y guarda precio y existencia por tienda. \texttt{detalle\_venta} resuelve venta--producto y conserva el precio histórico.

\section{Modelos de datos}
\subsection{Conceptual}
El modelo conceptual muestra entidades y cardinalidades sin decisiones físicas. Una venta pertenece a una tienda, cliente, empleado y estado; contiene uno o más detalles y admite uno o más pagos según el negocio.
\modelo{../modelos/modelo_conceptual.png}{Modelo conceptual entidad--relación.}
\subsection{Lógico}
El modelo lógico incorpora atributos, identificadores, obligatoriedad y claves candidatas. El símbolo \# identifica clave primaria, * referencia y UQ unicidad.
\modelo{../modelos/modelo_logico.png}{Modelo lógico normalizado.}
\subsection{Relacional}
El modelo relacional fija tablas, PK, FK y tipos Oracle. La FK compuesta de \texttt{venta(id\_empleado,id\_tienda)} garantiza declarativamente que el empleado pertenece a la tienda de la venta.
\modelo{../modelos/modelo_relacional.png}{Modelo relacional Oracle.}

\section{Normalización}
\textbf{Primera forma normal (1FN).} Cada columna contiene un valor atómico. Productos y pagos repetibles no se almacenan en columnas de venta, sino en \texttt{detalle\_venta} y \texttt{pago}. Direcciones, teléfonos e identificaciones se conservan como valores simples según el alcance.

\textbf{Segunda forma normal (2FN).} Las tablas con PK compuesta dependen de la clave completa: en \texttt{catalogo\_producto}, precio y existencia corresponden al par tienda--producto; en \texttt{detalle\_venta}, cantidad, precio aplicado y subtotal corresponden al par venta--producto. No existen dependencias parciales.

\textbf{Tercera forma normal (3FN).} Los atributos descriptivos de país, departamento, municipio, cargo, categoría, marca, estados y métodos están en sus catálogos. Los datos personales no se repiten entre empleado y cliente. La categoría y marca dependen del producto; la ubicación se deriva mediante sus FK; por tanto, ningún atributo no clave depende transitivamente de otra columna no clave dentro de su relación.

\section{Integridad y supuestos}
Se definieron PK en todas las entidades, FK para cada relación, \texttt{NOT NULL} en atributos obligatorios, unicidad para nombres de catálogos, correo, identificación y claves naturales compuestas, y \texttt{CHECK} para importes/cantidades positivas, existencia no negativa, estados permitidos y subtotal calculado.

Supuestos adoptados:
\begin{itemize}
\item Precio y existencia son específicos de tienda, porque el dataset los ubica en \texttt{CATALOGO\_PRODUCTO}; el producto maestro no los duplica.
\item Los 45 correos ausentes pertenecen a clientes y se cargan como \texttt{NULL}. El correo de cada empleado sí está presente y su unicidad se hereda de \texttt{persona}.
\item Una venta \texttt{REGISTRADA} puede tener abonos parciales. El archivo contiene 117 casos; la consulta 8 muestra total pagado y saldo.
\item No se usan triggers ni rutinas. Las reglas agregadas (venta con detalle, conciliación de pago) y la fecha relativa a \texttt{SYSDATE} se verifican con consultas, como permite el enunciado.
\item El precio histórico del detalle puede diferir del vigente sin constituir inconsistencia.
\end{itemize}

\section{Carga y conciliación}
La carga fue generada de manera determinista desde el Excel: textos se escapan, números usan punto decimal, fechas se expresan con \texttt{TO\_DATE(...,'YYYY-MM-DD')} y valores ausentes con \texttt{NULL}. Se ejecuta en orden de dependencias y termina en \texttt{COMMIT}.

\begin{longtable}{lr@{\qquad}lr}\toprule
Tabla&Filas&Tabla&Filas\\\midrule\endhead
PAIS&3&DEPARTAMENTO&32\\MUNICIPIO&80&TIPO\_TIENDA&4\\TIPO\_IDENTIFICACION&3&CARGO&4\\CATEGORIA&10&MARCA&15\\ESTADO\_VENTA&3&METODO\_PAGO&4\\PERSONA&388&TIENDA&20\\EMPLEADO&88&CLIENTE&300\\PRODUCTO&114&CATALOGO\_PRODUCTO&1,424\\VENTA&1,500&DETALLE\_VENTA&3,687\\PAGO&1,389&\textbf{TOTAL}&\textbf{9,068}\\\bottomrule\end{longtable}

La auditoría reproducible del libro obtuvo cero claves compuestas duplicadas, cero detalles sin catálogo, cero subtotales incorrectos, cero ventas sin detalle, cero cruces empleado--tienda inválidos y cero ventas pagadas desbalanceadas. El archivo \texttt{validacion\_dataset.txt} conserva la evidencia y \texttt{tools/validar\_dataset.py} permite repetirla.

\section{Orden de ejecución en Oracle SQL Developer}
\begin{enumerate}
\item Conectarse a un esquema vacío de Oracle Database.
\item Abrir \texttt{03\_creacion\_DB.sql} y usar \emph{Run Script} (F5). Debe indicar 19 tablas creadas.
\item Abrir \texttt{03\_carga\_datos.sql} y usar F5. Revisar los 19 conteos y el total 9,068.
\item Ejecutar individualmente \texttt{04\_Consulta1.sql} a \texttt{04\_Consulta8.sql}.
\item Ejecutar \texttt{05\_validaciones.sql}; cada uno de los cuatro conteos debe devolver 0.
\end{enumerate}

Los scripts incluyen \texttt{WHENEVER SQLERROR EXIT SQL.SQLCODE ROLLBACK} o comentarios de propósito. La carga incluye \texttt{SET DEFINE OFF} para evitar que un carácter ampersand dentro de datos se interprete como variable de sustitución.

\section{Consultas entregadas}
\begin{longtable}{p{.6cm}p{4.1cm}p{10.5cm}}\toprule N.&Reporte&Decisión técnica\\\midrule\endhead
1&Ventas por tienda y ubicación&Preagrupa detalle por venta; excluye anuladas y conserva tiendas sin ventas.\\
2&Ventas por tipo de tienda&Cuenta tiendas distintas y ventas sin multiplicarlas por detalle.\\
3&Productos más vendidos&Suma unidades y subtotal de ventas no anuladas.\\
4&Desempeño de empleados&Incluye empleados sin ventas mediante \texttt{LEFT JOIN}; excluye anuladas.\\
5&Clientes con mayor compra&Incluye exclusivamente ventas con estado \texttt{PAGADA}.\\
6&Facturación por categoría y marca&Agrupa los detalles no anulados por ambas dimensiones.\\
7&Uso de métodos de pago&Cuenta pagos reales y montos recibidos, excluyendo anuladas.\\
8&Ventas pendientes&Preagrupa detalles y pagos por separado, usa \texttt{LEFT JOIN} y calcula diferencia para \texttt{REGISTRADA}.\\\bottomrule\end{longtable}

El patrón de preagregación evita el producto cartesiano que se produciría al unir simultáneamente varios detalles y varios pagos de una misma venta.

\section{Consultas de validación}
\texttt{05\_validaciones.sql} comprueba: (1) empleado y tienda de venta coincidentes; (2) pagos iguales al total en toda venta pagada; (3) inexistencia de ventas sin detalle; y (4) contrataciones futuras. Las cuatro deben devolver cero. Las primeras dos satisfacen explícitamente el mínimo de la rúbrica y las otras dos amplían la cobertura.

\section{Herramientas y archivos fuente}
\begin{itemize}
\item Oracle Database y Oracle SQL Developer para DDL, DML, DQL y evidencia final.
\item Oracle SQL Developer Data Modeler para el modelo fuente; el archivo \texttt{.dmd} se entrega junto a su carpeta asociada.
\item Python 3 y openpyxl para generar y auditar la carga sin transcripción manual.
\item Graphviz para exportaciones SVG/PNG incorporadas en este manual.
\item XeLaTeX para producir los documentos PDF.
\end{itemize}

\section{Lista de comprobación para entrega}
Antes de subir el ZIP: abrir ambos PDF; abrir el \texttt{.dmd} sin separar su carpeta; probar DDL y carga desde un esquema vacío; confirmar 9,068 filas; ejecutar los ocho reportes; confirmar cuatro ceros; extraer una copia del ZIP y repetir el orden descrito.

\section{Limitación de la evidencia incluida}
La validación incluida se ejecutó directamente sobre el Excel y los artefactos se revisaron de forma estática en este entorno. La captura de una sesión Oracle autenticada debe obtenerse en la conexión oficial del estudiante al efectuar la prueba final; no se fabricaron capturas de una ejecución inexistente.
\end{document}
